% slides/02_motivation.tex


\begin{frame}{¿Por qué EEG cuantitativo?}
  \begin{columns}[T]
    \begin{column}{0.55\textwidth}
      \begin{itemize}
        \item El análisis \emph{clásico} del EEG parte de la \alert{inspección visual} — subjetiva, lenta, dependiente del experto
        \item[]
        \item El qEEG extrae \alert{rasgos numéricos objetivos y reproducibles} a partir de las señales crudas
        \item[]
        \item Posibilita los estudios poblacionales de gran escala, las interfaces cerebro-computadora en tiempo real, y el descubrimiento de biomarcadores
      \end{itemize}
    \end{column}
    \begin{column}{0.42\textwidth}
    	Confluencia de:
			\begin{itemize}
				\item Sistemas de alta densidad (256+ electrodos)
				\item Pipelines extendibles de software libre (MNE, BCI2000)
				\item Herramientas modernas de aprendizaje maquinal
			\end{itemize}
      \vspace{0.5em}
      \begin{block}{Referencias clave}
        \small
        Nuwer \textit{et al.}\ (1998) establecen estándares básicos de qEEG \cite{Nuwer1999} \\
        \vspace{5pt}
        Makeig \textit{et al.}\ (2004) consolidan el análisis de la dinámica cerebral \cite{Makeig2004}.
      \end{block}
    \end{column}
  \end{columns}
\end{frame}
